% !TEX program = lualatex
%==============================================================================
%  Adjoint-State Gradients on the Stratified Voxel Grid.
%  Waveform-inversion gradients through the Foldy--Lax solve, and their
%  Dietrich--Kormendi limit.
%
%  Self-contained LuaLaTeX source. Compile with:
%      latexmk -lualatex AdjointStateGradients.tex
%==============================================================================
\documentclass[11pt,a4paper]{article}

% --- LuaLaTeX font handling -------------------------------------------------
\usepackage{fontspec}
\setmainfont{Latin Modern Roman}

% --- Mathematics ------------------------------------------------------------
\usepackage{amsmath}
\usepackage{amssymb}
\usepackage{mathtools}

% --- Layout and typography --------------------------------------------------
\usepackage[margin=1in]{geometry}
\usepackage{microtype}
\usepackage{xcolor}
\usepackage{booktabs}
\usepackage{array}
\newcolumntype{L}[1]{>{\raggedright\arraybackslash}p{#1}}

% --- Hyperlinks -------------------------------------------------------------
\usepackage[colorlinks=true,
            linkcolor=blue!55!black,
            citecolor=blue!55!black,
            urlcolor=blue!55!black]{hyperref}
\hypersetup{
  pdftitle={Adjoint-State Gradients on the Stratified Voxel Grid},
  pdfauthor={Tod Nestor}
}
% Let long code paths break after underscores and dots.
\makeatletter
\g@addto@macro\UrlBreaks{\do\_\do\.}
\makeatother

% --- Macros -----------------------------------------------------------------
\newcommand{\ii}{\mathrm{i}}
\newcommand{\rd}{\mathrm{d}}
\newcommand{\Tz}{T_0}
\newcommand{\Gz}{G_0}
\newcommand{\bW}{\mathbf{W}}
\newcommand{\bR}{\mathbf{R}}
\newcommand{\bI}{\mathbf{I}}
\newcommand{\T}{\mathsf{T}}
\newcommand{\unval}{\textcolor{red!60!black}{\textbf{unvalidated}}}
\newcommand{\meas}{\textcolor{green!40!black}{\textbf{measured}}}

% Evidence files cited in this note.  Plain names, for searching:
% gate_t0_reciprocity  gate_9x9_source_convention  gate_sweep_dressed_equivalence
% gate_sweep_rung5_gmres  gate_dg0_absolute_magnitude  composed_matvec
\newcommand{\gtz}{\nolinkurl{g:t0_reciprocity}}

\title{Adjoint-State Gradients on the Stratified Voxel Grid\\[4pt]
\large Waveform-inversion gradients through the Foldy--Lax solve,\\
and their Dietrich--Kormendi limit}
\author{Tod Nestor}
\date{25 September 2026 \quad(design note; see \S\ref{sec:status} for what is validated)}

\begin{document}
\maketitle

\begin{abstract}
\noindent
The two-potential formulation of this programme --- a depth-stratified reference solved exactly,
plus a lateral contrast of cubic voxels summed by a Foldy--Lax solve carried by the reference
Green's tensor --- is the natural forward model for waveform inversion. Its first-order term
\emph{is} the Dietrich--Kormendi (1990) Born perturbation of the plane-wave reflectivity; the full
resolvent is its multiple-scattering generalisation. This note derives the adjoint-state gradient
of a waveform misfit with respect to every voxel's contrast. The gradient costs one extra Krylov
solve per source, frequency and transverse wavenumber, independent of the number of voxels, and
the adjoint system is the \emph{forward} system at $-k_y$, conjugated by $\mathbf{W}_{\pm} = \mathbf{I}_{\pm}\bW$,
the Voigt weight times a displacement--strain parity --- the nine-component form of the thesis
identity $\mathcal{H}^\dagger(k_y) = \mathcal{H}(-k_y)$. That identity requires matching
reciprocity laws of both $\Gz$ and $\Tz$, and both are measured here. The $\Gz$ law holds on the
assembled sweep operator to $10^{-13}$ or better, in a whole space and in a stratified
background, and the Foldy--Lax identity to $10^{-16}$ (gate A1). The $\Tz$ law is exact for the
Rayleigh cube carried by the sweep solver, and \emph{failed} by the resonance composite, which is
incidence-dependent and must not be used as it stands. A1 also corrected this note's first draft,
which stated the law with $\bW$ alone and the Foldy--Lax operator in the wrong order. The gradient
is implemented (\nolinkurl{sweep_gradient.py}) and passes a six-rung validation ladder
(\S\ref{sec:gates}). That covers the receiver transpose and the $\Tz$ derivative (complex step is
invalid, since $\Tz$ is not holomorphic in its contrast), an absolute check of the Born limit
against the exact thin-layer reflection, a dot-product test, and a Taylor test against the
nonlinear misfit. The last two reach spectral radius $0.74$--$0.80$.
\end{abstract}

\bigskip
\noindent\textbf{Companion series.}\enspace{\small\itshape One of a companion series on elastic
multiple scattering from cubic heterogeneities on a depth-stratified background, the T-matrix
programme $T=T_0(I-G_0T_0)^{-1}$: \textbf{(I)}~the stratified inter-site propagator
(\texttt{EwaldIntraPlanePropagator}); \textbf{(II)}~the contrast-source volume integral equation
and the seismic grid T-matrix lineage (\texttt{ContrastSourceVIE}); \textbf{(III)}~the directional
partial-wave sweep matvec and the Krylov multiple-scattering solve
(\texttt{DirectionalSweepSolver}); \textbf{(IV)}~adjoint-state gradients through that solve, for
waveform inversion (\texttt{AdjointStateGradients}). This is note~(IV).}

\tableofcontents
\bigskip

%==============================================================================
\section{What already exists}\label{sec:exists}
%==============================================================================

Nothing below is new machinery until \S\ref{sec:voxgrad}. The inversion problem sits on four
things this programme and its sibling repository have already built, and one thing the thesis
already proved.

\begin{center}
\small
\begin{tabular}{@{}L{0.30\textwidth}L{0.34\textwidth}L{0.28\textwidth}@{}}
\toprule
Object & Where & Validated to \\
\midrule
Forward operator $(\bI-\Gz\Tz)$, 2\textonehalf-D, stratified $\Gz$
  & note (III); \texttt{sweep\_solver}, \texttt{directional\_sweeps}
  & GMRES $=$ LU $=$ Neumann $1.2\times10^{-12}$; $\Delta\Gz$ vs Kennett $2\times10^{-10}$ \\
Forward operator, 3-D
  & note (III) \S``three-dimensional case'': real-space tables, reverberation spectral
  & designed and costed ($2.1$\,h build at $16^3$); not gated end to end \\
Composed $\mathsf{P}^z+\mathsf{P}^x$ on the slab kernel
  & \texttt{scripts/composed\_matvec.py}
  & \emph{blocked}: GATE 1 fails, lattice aliasing. Not needed here \\
Adjoint of the Foldy--Lax operator
  & thesis \texttt{VariationalSum.tex}, Eqs.~\texttt{sym}, \texttt{HsymDef}, \texttt{Htdef}
  & proved, in the thesis's own representation \\
1-D Fréchet derivatives, $=$ Dietrich--Kormendi
  & \nolinkurl{SeismicInversion/latex/frechet_derivatives.tex} \S``Equivalence to D\&K''
  & AD vs tangent-linear $<10^{-10}$; vs finite differences $<10^{-5}$ \\
1-D Levenberg--Marquardt, Hessian hierarchy
  & same, \S\S``Hessian'', ``Global Matrix Method'', ``LM inversion''
  & gated in \texttt{Kennett\_Reflectivity}, \texttt{GlobalMatrix} \\
\bottomrule
\end{tabular}
\end{center}

\noindent Two consequences shape the rest of this note. First, the adjoint is built on the
\emph{sweep} operator of note (III), not on the composed slab matvec, whose blocker is a
representation problem the sweeps do not have. Second, the reference-layer derivatives are a
solved problem in one dimension; what is missing is the gradient with respect to the
\emph{voxel} contrasts, through the multiple-scattering solve. That is the subject here.

%==============================================================================
\section{Forward model and misfit}\label{sec:forward}
%==============================================================================

Fix a source $s$, frequency $\omega$ and, in 2\textonehalf-D, a transverse wavenumber node
$k_y$. The nine-component exciting field $\psi$ on the voxel grid solves
\begin{equation}\label{eq:fl}
  (\bI - \Gz\Tz)\,\psi = \psi^{\mathrm{inc}},
\end{equation}
with $\Tz = \operatorname{blockdiag}_i \Tz^{(i)}(m_i)$ the per-voxel cube T-matrices and
$m_i = (\Delta\lambda_i,\Delta\mu_i,\Delta\rho_i)$ the contrast of voxel $i$ from the local
reference. The data are the reference response plus the field radiated by the contrast sources
$\tau = \Tz\psi$,
\begin{equation}\label{eq:data}
  d = d^{\mathrm{ref}} + \bR\,\Tz\,\psi ,
\end{equation}
where $\bR$ carries the stratified Green's tensor from each voxel to the receivers (the
$\mathcal{P}^{z,R}_{s}$ of the thesis \texttt{LSstrat}). With residual $r = d - d^{\mathrm{obs}}$,
\begin{equation}\label{eq:misfit}
  \mathsf{M} = \tfrac12 \sum_{s,\omega,k_y} w_{k_y}\, r^{\mathsf H} r .
\end{equation}
Moduli are complex (constant-$Q$), so every adjoint below is a \emph{transpose}, never a
Hermitian conjugate: the operators are complex-symmetric in a weighted sense, not Hermitian. The
only conjugate in the problem is the one on the residual.

%==============================================================================
\section{The voxel gradient}\label{sec:voxgrad}
%==============================================================================

Perturb $\Tz\to\Tz+\delta\Tz$. From \eqref{eq:fl},
$\delta\psi = (\bI-\Gz\Tz)^{-1}\Gz\,\delta\Tz\,\psi$, and so
\begin{equation}
  \delta d = \bR\bigl[\bI + \Tz(\bI-\Gz\Tz)^{-1}\Gz\bigr]\delta\Tz\,\psi
           = \bR\,(\bI-\Tz\Gz)^{-1}\,\delta\Tz\,\psi ,
\end{equation}
by the push-through identity $\bI + A(\bI-BA)^{-1}B = (\bI-AB)^{-1}$. Define the adjoint field
$\tilde\psi$ by
\begin{equation}\label{eq:adj}
  \boxed{\;(\bI - \Gz^{\T}\Tz^{\T})\,\tilde\psi = \tilde{\mathbf b},\qquad \tilde{\mathbf b} = \bR^{\T}\,\bar r\;}
\end{equation}
so that $\bar r^{\T}\bR(\bI-\Tz\Gz)^{-1} = \tilde\psi^{\T}$. The tilde follows the thesis
(\texttt{VariationalSum.tex}, Eqs.~\texttt{bvardef} and \texttt{btdef}). There the incident field
$\mathbf{b}$ is paired with $\tilde{\mathbf{b}} = [\mathcal{P}^{z,R}_{s}]^{\T}\tilde{\mathbf{f}}$,
radiated from the receiver with directivity $\tilde{\mathbf{f}}$. The right-hand side of
\eqref{eq:adj} is that field with $\tilde{\mathbf{f}} = \bar r$, and $\tilde\psi$ is its exciting
field. Then
$\delta\mathsf{M} = \operatorname{Re}\sum \tilde\psi^{\T}\delta\Tz\,\psi$, and because $\Tz$ is block-diagonal
the gradient is a local $9\times9$ contraction at each voxel:
\begin{equation}\label{eq:grad}
  \boxed{\;
  \frac{\partial\mathsf{M}}{\partial m_{i,\ell}}
  = \operatorname{Re}\sum_{s,\omega,k_y} w_{k_y}\;
    \tilde\psi_i^{\T}\,\frac{\partial\Tz^{(i)}}{\partial m_{i,\ell}}\,\psi_i ,
  \qquad \ell\in\{\Delta\lambda,\Delta\mu,\Delta\rho\}. \;}
\end{equation}
The right-hand side of \eqref{eq:adj} is the residual injected at the receivers and carried to the
voxels by the transposed receiver operator. By reciprocity $\bR^{\T}$ is the reference field
radiated \emph{from} the receivers, with the same weight $\mathbf{W}_{\pm}$ as \S\ref{sec:selfadj}:
$G(r\!\leftarrow\!v)^{\T} = \mathbf{W}_{\pm}\,G(v\!\leftarrow\!r)\,\mathbf{W}_{\pm}^{-1}$. For a receiver
recording displacement, $\bR = \Pi_u\,G(r\!\leftarrow\!v)$ with $\Pi_u$ selecting the three displacement
rows, and $\mathbf{W}_{\pm}^{-1}\Pi_u^{\T} = \Pi_u^{\T}$, so
\begin{equation}\label{eq:rt}
  \bR^{\T}\bar r = \mathbf{W}_{\pm}\cdot\bigl[\text{field at the voxels from a point \emph{force}
  $\bar r$ at the receiver}\bigr],
\end{equation}
which the existing layered incident-field code already computes. Both statements are measured in
\S\ref{sec:a1b}. In 2\textonehalf-D
the observed data live in real $y$, so the residual enters node $k_y$ with the phase and weight
of the $k_y$ quadrature, $w_{k_y}e^{\ii k_y y_r}\bar r(y_r)$.

\paragraph{Status.} Equations \eqref{eq:adj}--\eqref{eq:grad} are an exact algebraic identity for
any $\Gz$ and $\Tz$; nothing about the physics is assumed. Their implementation passes the
dot-product and Taylor tests A4 and A5 (\S\ref{sec:a4}, \S\ref{sec:a5}).

\paragraph{The derivative of the cube T-matrix.} $\partial\Tz/\partial m$ is not the Born weight.
The Rayleigh cube carries the \emph{effective} contrasts
$(\Delta\rho^*,\Delta\lambda^*,\Delta\mu^*_{\mathrm{diag}},\Delta\mu^*_{\mathrm{off}})$, which
are nonlinear in $(\Delta\lambda,\Delta\mu,\Delta\rho)$ through the Eshelby amplification factors
(\nolinkurl{effective_contrasts._compute_amplification_factors}). That is the point: the
gradient is exact for the resummed intra-voxel scattering, which a Born kernel is not.

The first draft of this note proposed a complex-step derivative. That is invalid here. Complex
step, and its generalisation the Cauchy contour integral, require $\Tz$ to be
\emph{holomorphic} in $m$, and it is not: \nolinkurl{compute_cube_tmatrix} applies its real form
factors to the \emph{real parts} of the effective contrasts only
($\Delta\mu^* \mapsto \operatorname{Re}\Delta\mu^*\,f_\mu + \ii\operatorname{Im}\Delta\mu^*$, and
likewise for $\kappa$ and $\rho$). Every contrast-dependent ingredient is otherwise a smooth
rational function, so $\Tz$ is smooth as a map of \emph{real} $m$, away from the resonance poles of
the amplification factors. The derivative is therefore taken by Richardson-extrapolated central
differences along the real axis. There are three parameters per voxel type, so it costs nine $\Tz$
evaluations. At zero contrast the amplification factors are exactly 1, and the derivative
reduces to the Born weights $\omega^2V\,\bI_3\oplus0$ and $0\oplus V\,\partial c_{\mathrm{Voigt}}$
only as $ka\to0$. At finite $ka$ it carries the form factor, which the Born weight does not.
Gate A2 measures all of this (\S\ref{sec:a2}).

\subsection{The derivative, measured (gate A2)}\label{sec:a2}

\nolinkurl{gate_a2_t0_derivative} uses the project's validated test parameters ($\alpha=5$,
$\beta=3$, $\rho=2.5$; contrast $(2\ \text{GPa}, 1\ \text{GPa}, 0.1\ \text{g/cm}^3)$; $a=0.5$ km)
at $ka_S = 0.05$ and $0.3$, and at a complex frequency ($\omega(1+0.03\ii)$) as in the sweep
gates.

\begin{center}
\small
\setlength{\tabcolsep}{4pt}\begin{tabular}{@{}lccc@{}}
\toprule
Check & $ka_S=0.05$ & $ka_S=0.3$ & $ka_S=0.3$, complex $\omega$ \\
\midrule
Richardson error estimate, worst parameter & $2.5\times10^{-11}$ & $3.4\times10^{-13}$ & $2.9\times10^{-13}$ \\
observed order of raw differences & 2.00 (see text) & 2.00 & 2.00 \\
random direction vs $\sum_\ell u_\ell\,\partial_\ell\Tz$ & $1.0\times10^{-13}$ & $1.0\times10^{-13}$ & $1.6\times10^{-13}$ \\
\emph{must fail:} imaginary-step derivative & $4$--$7\times10^{-4}$ & $1.5$--$2.4\times10^{-2}$ & $1.5$--$2.4\times10^{-2}$ \\
\bottomrule
\end{tabular}
\end{center}

The imaginary-step row is what complex step would have delivered. It is wrong by an amount that
grows like $(ka)^2$, which is at least $10^7$ times the finite-difference error. The density
channel at $ka_S = 0.05$ is linear in $\Delta\rho$ to rounding, so its raw differences agree to
rounding and have no truncation order to measure. The gate reports it as linear rather than
asserting an order. (The first run asserted a fixed $10^{-3}$ for the imaginary-step disagreement
and an order for every channel; both failed for those reasons, not for any defect in $\Tz$.)

\paragraph{Born limit.} At zero contrast, the relative difference from the Born weights is:

\begin{center}
\small
\begin{tabular}{@{}lcccc@{}}
\toprule
$ka_S$ & 0.40 & 0.20 & 0.10 & 0.05 \\
\midrule
$\partial_{\Delta\lambda}$, $\partial_{\Delta\mu}$, $\partial_{\Delta\rho}$ &
$1.90\times10^{-2}$ & $4.79\times10^{-3}$ & $1.20\times10^{-3}$ & $3.00\times10^{-4}$ \\
\bottomrule
\end{tabular}
\end{center}

\noindent The fitted slope is $2.00$, $1.99$, $1.99$, which is the $(ka)^2$ convergence the form
factor predicts. The residual is identical in all three channels at every $ka$, with a coefficient
of about $0.119$ in $(k_S a)^2$. That is consistent with a single contrast-independent geometric
form factor at zero contrast. This is an observation; the code path that makes the three equal has
not been traced.

%==============================================================================
\section{The adjoint is the forward solve at \texorpdfstring{$-k_y$}{-ky}}\label{sec:selfadj}
%==============================================================================

Equation \eqref{eq:adj} needs a solver for the transposed operator. It does not need a new one.
Both factors obey one diagonal similarity law (measured below),
\begin{equation}\label{eq:laws}
  \Gz(k_y)^{\T} = \mathbf{W}_{\pm}\,\Gz(-k_y)\,\mathbf{W}_{\pm}^{-1},
  \qquad
  \Tz^{\T} = \mathbf{W}_{\pm}\,\Tz\,\mathbf{W}_{\pm}^{-1},
  \qquad
  \mathbf{W}_{\pm} = \mathbf{I}_{\pm}\bW,
\end{equation}
\begin{equation*}
  \bW = \operatorname{diag}\bigl(1,1,1,\ 1,1,1,\ \tfrac12,\tfrac12,\tfrac12\bigr),
  \qquad
  \mathbf{I}_{\pm} = \operatorname{diag}\bigl(1,1,1,\ -1,-1,-1,\ -1,-1,-1\bigr).
\end{equation*}
$\bW$ is the Voigt engineering factor. $\mathbf{I}_{\pm}$ is the parity between the displacement and
strain halves of the state: the two off-diagonal blocks of the $9\times9$ carry one spatial
derivative, taken at the receiver in one and at the source in the other, so exchanging source and
receiver flips their sign. The adjoint system \eqref{eq:adj} is the transpose of
$(\bI - \Tz\Gz)$, \emph{not} of $(\bI-\Gz\Tz)$, and
$\Gz^{\T}\Tz^{\T} = \mathbf{W}_{\pm}\,\Gz(-k_y)\Tz\,\mathbf{W}_{\pm}^{-1}$, so
\begin{equation}\label{eq:selfadj}
  \boxed{\;
  \bigl(\bI - \Tz\Gz(k_y)\bigr)^{\T} = \mathbf{W}_{\pm}\bigl(\bI - \Gz(-k_y)\Tz\bigr)\mathbf{W}_{\pm}^{-1}
  \;}
\end{equation}
and \eqref{eq:adj} becomes a forward solve at $-k_y$ with the right-hand side
$\mathbf{W}_{\pm}^{-1}\tilde{\mathbf b}$, solved for $\tilde\psi_{-}$, followed by $\tilde\psi = \mathbf{W}_{\pm}\tilde\psi_{-}$. This is the nine-component statement of the
thesis result: with $\Delta\mathcal{C}_{\mathit{eff}}(k_y) = \Delta\mathcal{C}^{\T}_{\mathit{eff}}(-k_y)$
and $\mathcal{P}_{r,s}(k_y) = \mathcal{P}^{\T}_{r,s}(-k_y)$ (\texttt{VariationalSum.tex},
Eq.~\texttt{sym}), the Lippmann--Schwinger operator
$\mathcal{H}(k_y) = \bI - \mathcal{P}_{r,s}(k_y)\Delta\mathcal{C}_{\mathit{eff}}(k_y)$ has adjoint
$\mathcal{H}^\dagger(k_y) = \mathcal{H}(-k_y)$ (Eq.~\texttt{Htdef}). The thesis used that identity
to build Schwinger--Levine variational estimates; here it makes the gradient cost exactly one
forward solve. Two practical consequences follow. The $k_y$ grid of note (III) is symmetric about
zero, so the $-k_y$ operator is already in the cache; and the forward and adjoint solves share an
operator and can share a Krylov space. In 3-D, where $k_y$ is not a parameter, the same law is a
pair swap on the real-space tables, and the adjoint reuses the tables built for the forward solve.

\paragraph{The $\Tz$ half, measured.} \gtz\ checks $\Tz\bW^{-1}$ for symmetry, i.e.\
$\Tz^{\T} = \bW\Tz\bW^{-1}$. A cube $\Tz$ has no displacement--strain block, so $\mathbf{I}_{\pm}$
commutes with it, and this is the second law of \eqref{eq:laws}:

\begin{center}
\small
\setlength{\tabcolsep}{4pt}\begin{tabular}{@{}lccc@{}}
\toprule
$\Tz$ & $ka_S$ & asymmetry of $\Tz\bW^{-1}$ & normal--shear block \\
\midrule
Rayleigh cube & 0.05, 0.1, 0.3 & $0$ & $0$ \\
Resonance composite, $n_{\mathrm{sub}}=3$, $\hat k = \hat z$ & 0.5 & $1.6\times10^{-2}$ & $2\times10^{-19}$ \\
Resonance composite, $n_{\mathrm{sub}}=3$, $\hat k$ oblique & 1.2 & $1.7\times10^{-1}$ & $4\times10^{-4}$ \\
\bottomrule
\end{tabular}
\end{center}

The Rayleigh result is structural and was expected: the cube is block-diagonal, the density
block is a multiple of $\bI_3$, and a cubic Voigt stiffness has no normal--shear coupling, so
$\bW$ commutes with $\Tz$ and the law reduces to plain symmetry. The gate's value is as a
regression guard, with the composite rows as its calibration leg: they show the check can fail.

The composite fails, and not by a rounding amount. It is
$\sum_n T_{\mathrm{loc}}\,\psi^{\mathrm{exc}}_n$ for sub-cell exciting fields driven by one
incident plane wave, so it carries the incident phase on its input side and none on its output
side, and it changes with $\hat k$. It is a response to one illumination, not a field-independent
operator, and \emph{no} such object can be reciprocal. The adjoint field arrives from the
receivers, from a different direction than the source field, so a gradient built on the composite
would be wrong by the asymmetry in the table. Above the Rayleigh range the options are to keep the
sub-cells as explicit unknowns in the global solve (so the local operator is $T_{\mathrm{loc}}$,
which is Rayleigh and reciprocal), or to rebuild the composite with a symmetric projection. The
27-component Galerkin $\Tz$ has not been checked.

\paragraph{The $\Gz$ half, measured on the assembled operator (gate A1).} The pieces had been
checked separately (the whole-space $\bW$-symmetry of \texttt{StratifiedComposition}; the $J_6$
law of the reverberation, note (III) \texttt{[D3]}), but not the assembled discrete sweep operator,
which is what \eqref{eq:selfadj} applies. \nolinkurl{gate_a1_sweep_reciprocity} materialises it
densely at $\pm k_y$, split into the lateral sweep, the same-plane reverberation and the
inter-plane coupling. Residuals are $\lVert A - B\rVert/\lVert A\rVert$ at the finer of two
quadratures; every one is flat under refinement.

\begin{center}
\small
\setlength{\tabcolsep}{4pt}\begin{tabular}{@{}lcc@{}}
\toprule
Claim & whole space & stratified \\
\midrule
$\Gz$ law \eqref{eq:laws}, lateral sweep & $2.8\times10^{-14}$ & $1.6\times10^{-16}$ \\
\quad same-plane reverberation & --- & $1.0\times10^{-13}$ \\
\quad inter-plane coupling & $1.0\times10^{-14}$ & $5.1\times10^{-16}$ \\
\quad \textbf{assembled} $\Gz$ & $2.1\times10^{-14}$ & $1.7\times10^{-15}$ \\
identity \eqref{eq:selfadj}, Rayleigh $\Tz$ & $4.4\times10^{-17}$ & $3.0\times10^{-16}$ \\
\midrule
\emph{must fail:} $\bW$ in place of $\mathbf{W}_{\pm}$ & $0.74$ & $0.29$ \\
\quad $\mathbf{W}_{\pm}$ without the $k_y$ flip & $0.41$ & $0.097$ \\
\quad $(\bI-\Gz\Tz)^{\T}$ for $(\bI-\Tz\Gz)^{\T}$ & $2.4\times10^{-3}$ & $0.095$ \\
\bottomrule
\end{tabular}
\end{center}

$\mathbf{W}_{\pm}$ was measured, not assumed. A free fit of
$[\Gz(k_y)^{\T}]_{ab}/[\Gz(-k_y)]_{ab}$ is an exact outer product $c_a/c_b$, spread
$1.4\times10^{-13}$ (whole space) and $2.6\times10^{-11}$ (stratified), with $c$ equal in both to the diagonal of $\mathbf{W}_{\pm}$ given in \eqref{eq:laws}. The whole space is at
$\omega/2\pi = 1+0.03\ii$, $k_y = 0.6$, pitch $0.25$; the stratified case is the fast-slab crust of
the sweep-solver tests, planes 7 and 11, $\omega/2\pi = 6$, $k_y=0.3$, $Q=2$; both $2\times4$
voxels. $k_y\neq0$ is essential, since at $k_y=0$ the SH channel decouples and the flip is
invisible.

\paragraph{Two errors this gate caught in the first draft of this note.} The law was stated with
$\bW$ alone. That transcribed the whole-space $\bW$-symmetry, which compares entries at the same
separation, as a transpose law, which reverses the separation; the reversal is where $\mathbf{I}_{\pm}$
comes from, and without it the law fails at $0.29$--$0.74$. The identity \eqref{eq:selfadj} was
also written for $(\bI-\Gz\Tz)^{\T}$; that operator does \emph{not} map to the forward operator
at $-k_y$ (third calibration row), and the adjoint system never needed it. Neither error would
have been visible without the dense measurement: both drafts are algebraically self-consistent.

\subsection{The receiver transpose, measured (gate A1b)}\label{sec:a1b}

\nolinkurl{gate_a1b_receiver_transpose} tests the law in 3-D, where there is no $k_y$ to flip and
it is a pure exchange of source and receiver. It checks the wavenumber-domain kernel
\nolinkurl{corrected_layered_9x9} first, then the real-space field of
\nolinkurl{layered_incident_field}. For the real-space check it builds the dense $\bR$ one voxel at
a time, with nine unit sources at each voxel read at the receiver, and compares it with the field
radiated \emph{from} the receiver. The model is the fast-slab crust of the sweep-solver tests,
with voxel planes 7 and 11. One receiver is in the background above them (the same material as the
voxels); the other is \emph{inside} the slab (a different material), which is the case where
taking the wrong interface's material once left an 83\% residual (\texttt{StratifiedComposition}).

\begin{center}
\small
\setlength{\tabcolsep}{4pt}\begin{tabular}{@{}lccc@{}}
\toprule
Check & same material & cross material & $\bW$ in place of $\mathbf{W}_{\pm}$ \\
\midrule
$k$-domain, $Q=2, 1000$, surface off/on & $\le6.5\times10^{-16}$ & $\le7.6\times10^{-16}$ & $0.30$--$0.40$ \\
real space, $k_r^{\max}=10/\text{pitch}$ (documented rule) & $1.2\times10^{-16}$ & $\mathbf{1.4\times10^{-2}}$ & $0.30$, $0.37$ \\
real space, $k_r^{\max}=20/\text{pitch}$ (asserted) & $9.6\times10^{-17}$ & $1.1\times10^{-9}$ & $0.30$, $0.37$ \\
real space, $k_r^{\max}=30/\text{pitch}$ & $9.6\times10^{-17}$ & $7.7\times10^{-12}$ & $0.30$, $0.37$ \\
form \eqref{eq:rt}, force at the receiver, $20/\text{pitch}$ & $1.8\times10^{-16}$ & $2.6\times10^{-10}$ & --- \\
\bottomrule
\end{tabular}
\end{center}

The node spacing is held at $\Delta k = 0.31$ while the cutoff grows. At $Q=2$ the free surface
changes nothing (it is damped out), so the free-surface leg is discriminating only at $Q=1000$,
which is why the $k$-domain check is run at both.

\paragraph{A defect in the transverse rule, outside the adjoint.} The law holds; the documented
rule does not resolve it. At $k_r^{\max} = 10/\text{pitch}$ the cross-material pair fails at
$1.4\times10^{-2}$, flat in the node spacing (in a first run on a $3\times2$ lattice, $1.53\times10^{-2}$ at 32 nodes and $1.52\times10^{-2}$ at 64). Raising
only the cutoff at fixed spacing brings it to $1.1\times10^{-9}$ and then $7.7\times10^{-12}$. So this
is truncation, not a convention error, and it appears only when the two planes are in different
materials. The construction subtracts the whole-space kernel \emph{of the receiver plane's medium}
spectrally and adds it back in closed form. When the media differ, what remains to be transformed
is not reverberation alone but also the difference between the two direct paths, and
\nolinkurl{measure_dg0_transverse_rule} fixed the $10/\text{pitch}$ cutoff on a single pair of
planes in the \emph{same} medium (planes 51 and 50 of the \texttt{gate\_sh\_impedance} model).
\texttt{layered\_stack\_table}, which builds the 3-D \emph{forward} tables, uses the same
construction and rule. Measured since (\nolinkurl{measure_stack_table_cross_material}), the defect
is there and larger than the ${\sim}1\%$ first guessed. With a fast slab half a pitch from the
planes, $10/\text{pitch}$ is off by up to $25\%$ of the operator's scale, and that includes the
reverberation block of a background plane next to the slab, not only cross-material pairs.
$20/\text{pitch}$ converges to ${\le}3\times10^{-4}$. A no-slab control is exact. The correction
is recorded in note (III). For the adjoint the remedy is simply the finer cutoff.

%==============================================================================
\section{The Dietrich--Kormendi limit}\label{sec:dk}
%==============================================================================

At zero contrast, $\psi = \psi^{\mathrm{inc}}$ and $\tilde\psi = \bR^{\T}\bar r$, the residual carried
into the reference by reciprocity. Equation \eqref{eq:grad} becomes
\begin{equation}\label{eq:dkgrad}
  \left.\frac{\partial\mathsf{M}}{\partial m_{i,\ell}}\right|_{\Tz=0}
  = \operatorname{Re}\sum_{s,\omega,k_y} w_{k_y}\;
    \bigl(\bR^{\T}\bar r\bigr)_i^{\T}\;
    \left.\frac{\partial\Tz^{(i)}}{\partial m_{i,\ell}}\right|_0\;
    \psi^{\mathrm{inc}}_i ,
\end{equation}
the reference Green's function to the receivers, times the Born vertex, times the reference field
from the source. That is the Dietrich--Kormendi Fréchet kernel,
$\delta R = \int G^{\uparrow}\,\delta\mathbf{P}\,G^{\downarrow}\,\rd z$, contracted with the
residual: a migration of the residual in the stratified reference, as Tarantola observed for the
homogeneous case. So the D\&K derivative is the first gradient of this inversion, and the full
formulation adds exactly two things: the multiple scattering between voxels, through
$(\bI-\Gz\Tz)^{-1}$, and the resummed intra-voxel scattering, through $\partial\Tz/\partial m$
away from zero.

This limit supplies the one \emph{absolute} gate of the ladder (A3). The earlier invariants and
self-consistency checks are all blind to an overall scale or sign; this one is not.

\subsection{The limit, measured (gate A3)}\label{sec:a3}

The first draft planned the AD Kennett Jacobian of \nolinkurl{Kennett_Reflectivity} as the arbiter.
That code's PP reflectivity is referenced to a water layer above a fluid--solid seabed, so matching
it would put the voxel side's sources and receivers in water, which the corrected layered kernel
does not support. The gate instead uses the D\&K kernel in its simplest exact form. In a uniform
medium the Born PP reflection of a layer of thickness $h$, referenced at its top, is exactly
\begin{equation}\label{eq:bornlayer}
  \delta R_{PP} = R_{\mathrm{lin}}(\delta m)\,\bigl(1 - e^{2\ii\omega\eta_P h}\bigr),
\end{equation}
with $R_{\mathrm{lin}}$ the Aki \& Richards linearised PP coefficient. That closed form shares no code
with the voxel side. On the voxel side, cubes of side $2a$ on a $2a\times2a$ lattice form a sheet at
mid-layer with Born source density $(2a)^{-2}\,\partial\Tz/\partial m|_0\,\psi^{\mathrm{inc}}$ per unit
area. This uses the gradient code's own Richardson derivative. The sheet radiates through
\nolinkurl{vertical_kernel_9x9}, the kernel the receiver maps and the vertical sweep carry, and is
projected onto up-going P by the mode bridge. P polarisations follow Aki \& Richards on both
sides. The sheet samples the layer at its midpoint and $\Tz$ carries its form factor, so the two
differ by $O(a^2)$. Any $O(1)$ defect (a sign, a missing $2\pi$ or $(2a)^{-2}$, a wrong vertex) would
leave a flat residual.

\begin{center}
\small
\setlength{\tabcolsep}{4pt}\begin{tabular}{@{}lccc@{}}
\toprule
Incidence & residual at $2a=0.4$ & residual at $2a=0.05$ km & slope in $a$ ($\Delta\lambda$, $\Delta\mu$, $\Delta\rho$) \\
\midrule
normal & $2.2\times10^{-2}$ & $3.3\times10^{-4}$ & 2.01, 2.01, 2.01 \\
$30^\circ$, in $k_x$ & $1.1\times10^{-2}$ & $1.7\times10^{-4}$ & 2.01, 2.01, 2.01 \\
$48.6^\circ$, split over $k_x$ and $k_y$ & $3.4\times10^{-3}$ & $4.1\times10^{-5}$ & 1.99, 2.11, 1.98 \\
\midrule
\emph{must fail:} vertex sign flipped & $2.0$ & $2.0$ & 0.00 \\
\quad $(2a)^{-2}$ omitted & $0.84$ & $1.0$ & $-0.08$ \\
\bottomrule
\end{tabular}
\end{center}

\noindent The medium is $\alpha=5$, $\beta=3$, $\rho=2.5$ at 1\,Hz; the largest residual per row is
quoted. The voxel Born vertex, the kernel's absolute normalisation and the polarisation conventions
are therefore right in absolute terms, not only consistent with one another. A3 does \emph{not}
cover the finite-lattice assembly or the adjoint chain (the unit tests of \S\ref{sec:impl} check
those against brute force). It also does not cover the lattice's non-specular orders, which are
evanescent for $2a$ below half a wavelength and are dropped, or a stratified reference.

%==============================================================================
\section{Reference-layer parameters}\label{sec:ref}
%==============================================================================

Note (III) found that the crust gradient belongs in the reference, not in the scattering series.
The reference parameters $m^{\mathrm{ref}}$ enter $d^{\mathrm{ref}}$, $\bR$, $\psi^{\mathrm{inc}}$
and $\Gz$. Differentiating \eqref{eq:fl}--\eqref{eq:data} and using
$\bR\Tz(\bI-\Gz\Tz)^{-1} = \bR(\bI-\Tz\Gz)^{-1}\Tz$,
\begin{equation}\label{eq:refgrad}
  \delta\mathsf{M} = \operatorname{Re}\sum \Bigl[
    \bar r^{\T}\delta d^{\mathrm{ref}}
    + \bar r^{\T}\delta\bR\,\tau
    + \tilde\psi^{\T}\Tz\,\delta\psi^{\mathrm{inc}}
    + (\Tz^{\T}\tilde\psi)^{\T}\,\delta\Gz\,\tau
  \Bigr],
  \qquad \tau = \Tz\psi .
\end{equation}
The first term is the 1-D problem, already solved and gated. The second and third are the
reference Green's function differentiated at the receivers and the source. The fourth contracts
the adjoint and forward contrast sources through $\partial\Gz/\partial m^{\mathrm{ref}}$ and costs
about one $\Gz$ matvec per reference parameter, which is acceptable because $\Gz$ depends only on
the reference. Whether \nolinkurl{GlobalMatrix.layered_greens} can be differentiated through the
existing \nolinkurl{gmm_torch} machinery has not been established. The recommended first scheme is
\emph{alternating}: update the reference with the existing 1-D machinery at fixed voxels, then the
voxels with \eqref{eq:grad} at fixed reference. Joint inversion waits for \eqref{eq:refgrad}.

%==============================================================================
\section{Cost, Gauss--Newton, and place in the lineage}\label{sec:cost}
%==============================================================================

Per source, frequency and $k_y$ node, the gradient costs one forward and one adjoint Krylov solve
on the same operator \eqref{eq:selfadj}, whatever the number of voxels or receivers. In the
frequency domain the stored state is $\psi$ alone, nine complex numbers per voxel, with none of the
wavefield checkpointing a time-domain adjoint needs. A Fr\'echet-derivative product
$\mathsf{F}\,\delta m = \bR(\bI-\Tz\Gz)^{-1}\delta\Tz\,\psi$ is one further solve, so Gauss--Newton by
inner conjugate gradients costs two solves per inner iteration, and the second-order Born Hessian
hierarchy of \nolinkurl{frechet_derivatives.tex} is the 1-D template for its approximations.

In the lineage of note (II), this is the distorted-Born iterative T-matrix inversion of Jakobsen
\& Ursin, with an analytic, resummed cell T-matrix and an exact stratified reference; and in the
thesis's own terms, the Schwinger--Levine functional already pairs the forward field with the
$-k_y$ field, which is the pairing \eqref{eq:grad} uses.

%==============================================================================
\section{Implementation}\label{sec:impl}
%==============================================================================

\nolinkurl{cubic_scattering/sweep_gradient.py} implements \eqref{eq:adj}--\eqref{eq:grad} for one
source, frequency and $k_y$ node on the 2\textonehalf-D directional-sweep solver. The sum over
$k_y$ nodes, with the quadrature weight and phase of \S\ref{sec:voxgrad}, is left to the caller.

\begin{center}
\small
\begin{tabular}{@{}L{0.30\textwidth}L{0.62\textwidth}@{}}
\toprule
Function & What it does, and the result it rests on \\
\midrule
\nolinkurl{misfit_and_gradient} & forward solve, data, misfit, then $\mathrm{Re}(\mathsf{F}^{\mathsf H}r)$
  through \nolinkurl{frechet_adjoint_product}; refuses caches that are not at $\pm k_y$ \\
\nolinkurl{frechet_vector_product} & the tangent-linear $\mathsf{F}\,\delta m = \bR(\bI-\Tz\Gz)^{-1}\delta\Tz\,\psi$
  by one further forward solve; the Gauss--Newton building block \\
\nolinkurl{frechet_adjoint_product} & $\mathrm{Re}(\mathsf{F}^{\mathsf H}q)$ for any data-space $q$:
  adjoint solve and local contraction \eqref{eq:grad} \\
\nolinkurl{adjoint_solve} & the forward solver at $-k_y$ on $\mathbf{W}_{\pm}^{-1}\tilde{\mathbf b}$, then $\tilde\psi=\mathbf{W}_{\pm}\tilde\psi_{-}$ (A1) \\
\nolinkurl{rayleigh_t0_and_derivative} & Rayleigh $\Tz$ per voxel in its plane's medium, and
  $\partial\Tz/\partial m$ by one Richardson step (A2) \\
\nolinkurl{whole_space_receiver_map}, \nolinkurl{layered_receiver_map} & a \emph{dense} $\bR$
  (one row per recorded component), so $\bR^{\T}$ is exact by construction; the force form
  \eqref{eq:rt} (A1b) is the route to take only if $\bR$ becomes too large to store \\
\bottomrule
\end{tabular}
\end{center}

\noindent Only the Rayleigh cube is supported; the resonance composite fails the $\Tz$ law. The
numerical settings (GMRES tolerance, iteration cap, Richardson step) are required arguments with
no defaults. Wiring the gradient to the $\omega$ annotations exposed that
\nolinkurl{compute_cube_tmatrix} and six helpers were annotated \texttt{omega: float} while every
sweep gate passes a complex $\omega$. The annotations were widened to \texttt{complex}; the
runtime is unchanged.

\paragraph{Tests} (\nolinkurl{cubic_scattering/tests/test_sweep_gradient.py}, fourteen):
\begin{itemize}
\item the adjoint solve against a dense solve of $(\bI-\Gz^{\T}\Tz^{\T})\tilde\psi=\tilde{\mathbf b}$, to $10^{-9}$;
\item the gradient against brute-force central differences of the misfit, to $10^{-6}$. In a
  whole space this is at weak and moderate contrast. In the fast-slab crust it uses stratified
  forward and adjoint caches, per-plane media, and a receiver \emph{inside} the slab;
\item the receiver map against the entries of the dense sweep operator $\Gz$, with asymmetric
  offsets, to $10^{-12}$;
\item on a uniform model, the layered receiver map reduces to the whole-space map to $10^{-9}$ at
  3\,Hz. At 0.3\,Hz the two differ by the physical seabed return, $9.7\times10^{-4}$, flat in the
  quadrature rule and falling with frequency; the first version of that test ran at 0.3\,Hz and
  failed for that reason;
\item $\mathsf{F}\,\delta m$ against central differences of the predicted data, to $10^{-7}$; the dot-product
  identity, to $10^{-10}$; and the misfit gradient equal to $\mathrm{Re}(\mathsf{F}^{\mathsf H}r)$;
\item $\partial\Tz/\partial m$ against an independent small-step difference; zero residual giving
  zero gradient; and three fail-fast diagnostics.
\end{itemize}
Deliberate defects were injected to show the tests can fail. Each of these turned at least one
test red: $\bW$ in place of $\mathbf{W}_{\pm}$; the adjoint solved at $+k_y$; the residual left
unconjugated; $\partial\Tz/\partial m$ frozen at its Born value; the lateral phase sign of $\bR$
flipped; and the multiple-scattering term dropped from $\mathsf{F}\,\delta m$.

These tests check the gradient of the \emph{discrete} misfit at moderate contrast. They do not
replace A3 (absolute agreement with the D\&K limit), A4, or A5 (the Taylor test).

\subsection{The dot product, measured (gate A4)}\label{sec:a4}

\nolinkurl{gate_a4_dot_product} checks $\mathrm{Re}(q^{\mathsf H}\mathsf{F}\,\delta m) =
\delta m\cdot\mathrm{Re}(\mathsf{F}^{\mathsf H}q)$ for five random pairs $(\delta m, q)$ per case. $\mathsf{F}\,\delta m$
is computed forward, and $\mathsf{F}^{\mathsf H}q$ by the adjoint route the gradient uses. The two share only
$\psi$ and $\partial\Tz/\partial m$, and there is no step, so the identity holds to solver
tolerance only if the adjoint really is the transpose. That covers the A1 and A1b laws, the
receiver transpose, the conjugate on $q$ and the real part, all at once. GMRES runs at
$10^{-12}$, and the spectral radius $\operatorname{spr}(\Gz\Tz)$ is measured for each case.

\begin{center}
\small
\setlength{\tabcolsep}{4pt}\begin{tabular}{@{}lccccc@{}}
\toprule
Case & contrast & $\operatorname{spr}(\Gz\Tz)$ & mismatch & \emph{must fail:} Born $\mathsf{F}$ & \emph{must fail:} $\bW$ for $\mathbf{W}_{\pm}$ \\
\midrule
whole space, $ka_S=0.26$ & $1\times$ & 0.011 & $4.1\times10^{-13}$ & $3.2\times10^{-2}$ & $0.11$ \\
 & $4\times$ & 0.041 & $4.4\times10^{-13}$ & $0.12$ & $0.70$ \\
 & $8\times$ & 0.075 & $2.4\times10^{-11}$ & $8.1$ & $3.4$ \\
stress, $ka_S=0.52$ & $25\times$ & 0.427 & $7.6\times10^{-12}$ & $0.22$ & $0.72$ \\
 & $40\times$ & 0.804 & $6.6\times10^{-12}$ & $1.1$ & $24$ \\
stratified, receiver in slab & $1\times$ & 0.056 & $1.8\times10^{-13}$ & $2.7\times10^{-2}$ & $5.0\times10^{-2}$ \\
 & $4\times$ & 0.196 & $1.1\times10^{-12}$ & $6.2\times10^{-2}$ & $0.15$ \\
 & $8\times$ & 0.334 & $3.9\times10^{-12}$ & $9.3\times10^{-2}$ & $0.23$ \\
\bottomrule
\end{tabular}
\end{center}

\noindent Contrast is a multiple of the project's test contrast $(2\ \text{GPa},1\ \text{GPa},
0.1\ \text{g/cm}^3)$, perturbed per voxel. The identity holds to $\le2.4\times10^{-11}$ everywhere,
including $\operatorname{spr}=0.80$, and both calibration legs fail in every row.

\paragraph{How strong the multiple scattering can get.} With \emph{physical} Rayleigh $\Tz$
inside its validity ($ka_S\le0.3$) and the project's contrasts, $\operatorname{spr}(\Gz\Tz)$ stays at about $0.3$ or
below. Reaching $\operatorname{spr}\approx0.8$ needs $ka_S\approx0.5$ and $\Delta\mu\approx1.8\mu_0$, outside the
Rayleigh cube's validity. The dot product is an algebraic identity, so the \emph{stress} rows are
run there, labelled as such, to exercise the resolvent. The $\operatorname{spr}=0.92$ of the rung-5 solver gate
used random $\Tz$ blocks, not physical ones, so the ladder's earlier plan to run A5 ``at the rung-5
spectral radius'' with physical $\Tz$ was not achievable as written.

\subsection{The Taylor test, measured (gate A5)}\label{sec:a5}

A4 shows the adjoint is the transpose of the tangent-linear model. The Taylor test shows that
model is the derivative of the \emph{misfit}. For a random direction $\delta m$,
\begin{equation}\label{eq:taylor}
  \mathrm{Rem}_2(\varepsilon) = \mathsf{M}(m+\varepsilon\,\delta m) - \mathsf{M}(m) - \varepsilon\,\nabla\mathsf{M}\cdot\delta m
  = O(\varepsilon^2)
\end{equation}
if and only if $\nabla\mathsf{M}$ is right. Any gradient error leaves an $O(\varepsilon)$ piece, and
$\mathrm{Rem}_2/\varepsilon^2$ must settle to $\tfrac12\,\delta m^{\T}(\nabla^2\mathsf{M})\,\delta m$.
\nolinkurl{gate_a5_taylor} evaluates $\mathsf{M}$ with the full nonlinear model at every $\varepsilon$:
$\Tz$ recomputed from the perturbed contrasts, the Foldy--Lax solve and the receiver map. It uses
data from a different ``true'' model, so the residual is $O(1)$. $\varepsilon$ runs from $10^{-1}$
to $10^{-1}/2^{10}$ of the base contrast.

\begin{center}
\small
\setlength{\tabcolsep}{4pt}\begin{tabular}{@{}lccccc@{}}
\toprule
Case & spr & slope, $\mathrm{Rem}_2$ & steady to & $\mathrm{Rem}_1$ & Born $\nabla\mathsf{M}$ \\
\midrule
whole space, $ka_S=0.26$, $8\times$ & 0.071 & 2.000 & $3.4\times10^{-6}$ & 0.993 & 1.021 \\
stratified, receiver in slab, $8\times$ & 0.358 & 2.004 & $7.5\times10^{-5}$ & 1.001 & 1.007 \\
stress, $ka_S=0.52$, $40\times$ & 0.739 & 2.012 & $1.9\times10^{-4}$ & 0.988 & 1.005 \\
\bottomrule
\end{tabular}
\end{center}

\noindent Columns: spr is $\operatorname{spr}(\Gz\Tz)$; ``steady to'' is the relative change of
$\mathrm{Rem}_2/\varepsilon^2$ over the finest three $\varepsilon$; the last two are the slopes of the
calibration legs, which must be $\approx1$. $\mathrm{Rem}_1$ omits the gradient term, and the
Born-only gradient back-propagates the residual without the adjoint solve. Their slopes are
asymptotic, fitted over the finest three $\varepsilon$. No solver floor appears: $\mathrm{Rem}_2/\varepsilon^2$ is steady to four digits
at $\varepsilon\approx10^{-4}$ with GMRES at $10^{-12}$.

Two design errors in the first run, neither a gradient defect, are recorded because they would
recur. The calibration slopes were first fitted over the whole range, and the gate reported itself
blind. Where $\nabla\mathsf{M}\cdot\delta m$ is small beside the curvature ($\mathrm{Rem}_1$ even changes sign), or the
Born error is small ($\operatorname{spr}\approx0.07$), the $O(\varepsilon^2)$ term dominates at large $\varepsilon$,
and only the small-$\varepsilon$ end shows the $O(\varepsilon)$ defect the legs exist to expose. In the
stress case $\nabla\mathsf{M}\cdot\delta m\approx-0.04$ against a curvature of $2.6$, so $\mathrm{Rem}_1$ turns linear
only for $\varepsilon\ll0.017$, and the first $\varepsilon$ grid, stopping at $10^{-1}/2^6$, did not
reach it.

%==============================================================================
\section{Validation ladder}\label{sec:gates}
%==============================================================================

In order. Each rung is a prerequisite for trusting the next. Following note (III), every invariant
is paired with at least one absolute comparison, because reciprocity and dot-product tests are
homogeneous of degree one and blind to an overall scale error.

\begin{center}
\small
\begin{tabular}{@{}lL{0.34\textwidth}L{0.34\textwidth}L{0.17\textwidth}@{}}
\toprule
& Claim & Arbiter & Status \\
\midrule
A0 & $\Tz^{\T} = \bW\Tz\bW^{-1} = \mathbf{W}_{\pm}\Tz\mathbf{W}_{\pm}^{-1}$ & direct measurement, with failing calibration leg & \meas{} \gtz \\
A1 & $\Gz(k_y)^{\T} = \mathbf{W}_{\pm}\Gz(-k_y)\mathbf{W}_{\pm}^{-1}$ on the assembled sweep operator, and
     identity \eqref{eq:selfadj} & dense materialisation at $\pm k_y$, three failing calibration legs
     & \meas{} \nolinkurl{g:a1_sweep_reciprocity} \\
A1b & $G(r\!\leftarrow\!v)^{\T} = \mathbf{W}_{\pm}G(v\!\leftarrow\!r)\mathbf{W}_{\pm}^{-1}$ and form
      \eqref{eq:rt} of $\bR^{\T}$ & $k$-domain kernel; real-space dense $\bR$, two failing calibration legs
      & \meas{} \nolinkurl{g:a1b_receiver_transpose} \\
A2 & $\partial\Tz/\partial m$, and its Born limit & Richardson differences; directional consistency;
      Born weights as $ka\to0$; imaginary step must fail & \meas{} \nolinkurl{g:a2_t0_derivative} \\
A3 & D\&K limit \eqref{eq:dkgrad} & exact Born thin layer \eqref{eq:bornlayer} (absolute; slope 2 in
      pitch; two failing calibration legs) & \meas{} \nolinkurl{g:a3_dk_limit} \\
A4 & $\mathrm{Re}(q^{\mathsf H}\mathsf{F}\delta m) = \delta m\cdot\mathrm{Re}(\mathsf{F}^{\mathsf H}q)$ & dot-product
     test, $\mathsf{F}$ forward vs $\mathsf{F}^{\mathsf H}$ adjoint, to $\operatorname{spr}=0.80$; two failing calibration legs
     & \meas{} \nolinkurl{g:a4_dot_product} \\
A5 & gradient \eqref{eq:grad} at finite contrast & Taylor test \eqref{eq:taylor}, slope 2, to
     $\operatorname{spr}=0.74$ (labelled stress case); two calibration legs at slope 1 & \meas{} \nolinkurl{g:a5_taylor} \\
\bottomrule
\end{tabular}
\end{center}

Every rung has passed. The two that exercise the multiple scattering, A4 and A5, reach
$\operatorname{spr}\approx0.36$ with physical $\Tz$ and $\operatorname{spr}\approx0.74$--$0.80$ in a labelled stress case outside
the Rayleigh cube's validity (\S\ref{sec:a4}).

%==============================================================================
\section{Status}\label{sec:status}
%==============================================================================

\begin{center}
\small
\begin{tabular}{@{}L{0.42\textwidth}L{0.32\textwidth}L{0.18\textwidth}@{}}
\toprule
Item & Status & Evidence \\
\midrule
Gradient identity \eqref{eq:adj}--\eqref{eq:grad} & exact algebra & \S\ref{sec:voxgrad} \\
Adjoint $=$ forward at $-k_y$ \eqref{eq:selfadj} & \meas, $\le3\times10^{-16}$ & A1 \\
$\Tz$ half of \eqref{eq:laws}, Rayleigh cube & \meas, exact & \gtz \\
$\Tz$ half, resonance composite & \textbf{fails}, 1.6--17\% & \gtz \\
$\Gz$ half, assembled, whole space and stratified & \meas, $\le2\times10^{-14}$ & A1 \\
Receiver transpose $\bR^{\T}$, form \eqref{eq:rt} & \meas, $\le10^{-9}$ at $k_r^{\max}=20/\text{pitch}$ & A1b \\
Transverse rule $10/\text{pitch}$, cross-material planes & \textbf{inadequate}, $1.4\times10^{-2}$ & A1b \\
$\partial\Tz/\partial m$, Richardson differences & \meas, $\le2.5\times10^{-11}$ & A2 \\
$\Tz$ holomorphic in $m$ (complex step) & \textbf{no}; off by $4\times10^{-4}$--$2\times10^{-2}$ & A2 \\
D\&K limit, uniform reference & \meas, $O(a^2)$ to $\le3.3\times10^{-4}$ & A3 \\
Dot product, $\mathsf{F}$ forward vs $\mathsf{F}^{\mathsf H}$ adjoint & \meas, $\le2.4\times10^{-11}$ to $\operatorname{spr}=0.80$ & A4 \\
$\operatorname{spr}(\Gz\Tz)$, physical Rayleigh $\Tz$ within validity & $\lesssim0.36$ & A4, A5 \\
Taylor test of $\nabla\mathsf{M}$ against the nonlinear misfit & \meas, slope 2.000--2.012 to $\operatorname{spr}=0.74$ & A5 \\
Reference-layer gradient & 1-D term gated elsewhere; coupling terms \unval & \S\ref{sec:ref} \\
\midrule
Gradient code, 2\textonehalf-D, one $k_y$ node & built; passes A0--A5 & \S\ref{sec:impl} \\
\bottomrule
\end{tabular}
\end{center}

The validation ladder is complete. A0--A2 measure every ingredient of the gradient: the adjoint as
the forward solve at $-k_y$ conjugated by $\mathbf{W}_{\pm}$, the receiver transpose, and a Richardson
derivative of $\Tz$. The code (\S\ref{sec:impl}) assembles them. A3 confirms its Born limit
against exact physics in absolute terms. A4 shows the adjoint is the exact transpose of the
tangent-linear model, and A5 that this model is the derivative of the nonlinear misfit, both into
strong multiple scattering.

What the ladder does not cover is equally specific:
\begin{itemize}
\item the sum over $k_y$ nodes, which is the caller's (\S\ref{sec:voxgrad});
\item the 3-D solver, and reference-layer parameters beyond the 1-D term (\S\ref{sec:ref});
\item $\Tz$ other than the Rayleigh cube: the resonance composite fails the law, and the Galerkin
  $\Tz$ is untested;
\item physical $\Tz$ at $\operatorname{spr}$ above about $0.36$, which lies outside the Rayleigh cube's validity;
\item the stratified A3.
\end{itemize}
The A1 and A1b evidence is a few voxels at one frequency per background: small, but the laws are
pairwise identities, so lattice size does not enter them. A2 covers the Rayleigh cube only; the resonance composite is
excluded (\S\ref{sec:selfadj}), and the 27-component Galerkin $\Tz$ is untested.
A1b also found a defect outside the adjoint: the transverse rule is under-resolved for
planes near a material contrast, which reaches the \emph{forward} 3-D tables at up to $25\%$ of
the operator's scale (\S\ref{sec:a1b}).

\begin{thebibliography}{9}\small
\bibitem{DK90} Dietrich, M.\ \& Kormendi, F.\ (1990). Perturbation of the plane-wave reflectivity
  of a depth-dependent elastic medium by weak inhomogeneities. \emph{Geophys.\ J.\ Int.} 100,
  203--214.
\bibitem{KD91} Kormendi, F.\ \& Dietrich, M.\ (1991). Nonlinear waveform inversion of plane-wave
  seismograms in stratified elastic media. \emph{Geophysics} 56, 664--674.
\bibitem{Tar84} Tarantola, A.\ (1984). Inversion of seismic reflection data in the acoustic
  approximation. \emph{Geophysics} 49, 1259--1266.
\bibitem{Ple06} Plessix, R.-E.\ (2006). A review of the adjoint-state method for computing the
  gradient of a functional with geophysical applications. \emph{Geophys.\ J.\ Int.} 167, 495--503.
\bibitem{TTL05} Tromp, J., Tape, C.\ \& Liu, Q.\ (2005). Seismic tomography, adjoint methods, time
  reversal and banana-doughnut kernels. \emph{Geophys.\ J.\ Int.} 160, 195--216.
\bibitem{JU15} Jakobsen, M.\ \& Ursin, B.\ (2015). Full waveform inversion in the frequency domain
  using direct iterative T-matrix methods. \emph{J.\ Geophys.\ Eng.} 12, 400--418.
\bibitem{Ken83} Kennett, B.\,L.\,N.\ (1983). \emph{Seismic Wave Propagation in Stratified Media}.
  Cambridge University Press.
\end{thebibliography}

\vfill
\noindent\rule{\linewidth}{0.4pt}\\
{\footnotesize The measured values in this note are reproducible from
\nolinkurl{scripts/gate_t0_reciprocity.py}, \nolinkurl{scripts/gate_a1_sweep_reciprocity.py},
\nolinkurl{scripts/gate_a1b_receiver_transpose.py},
\nolinkurl{scripts/gate_a2_t0_derivative.py}, \nolinkurl{scripts/gate_a3_dk_limit.py},
\nolinkurl{scripts/gate_a4_dot_product.py} and \nolinkurl{scripts/gate_a5_taylor.py}.
Values quoted from other notes are reproducible from the gates those notes cite.}

\end{document}
